\documentclass[12pt,a4paper]{article}
\usepackage{amsmath,amssymb,amsthm}
\usepackage{geometry}
\usepackage{hyperref}
\usepackage{enumitem}
\geometry{margin=2.5cm}
\newtheorem{theorem}{Theorem}[section]
\newtheorem{lemma}[theorem]{Lemma}
\newtheorem{conjecture}[theorem]{Conjecture}
\newtheorem{proposition}[theorem]{Proposition}
\newtheorem{definition}[theorem]{Definition}
\newcommand{\R}{\mathbb{R}}
\newcommand{\C}{\mathbb{C}}
\newcommand{\Z}{\mathbb{Z}}
\newcommand{\T}{\mathbb{T}}
\newcommand{\calH}{\mathcal{H}}
\newcommand{\calD}{\mathcal{D}}
\newcommand{\tr}{\operatorname{Tr}}
\newcommand{\regdet}{\operatorname{det}_{\text{reg}}}
\newcommand{\SE}{\Xi}
\title{\bf The Generalised Riemann Helix and the Connes--Consani--Moscovici Spectral Triple:\\A Geometric Blueprint for the Hilbert--P\'olya Programme}
\author{}
\date{}

\begin{document}
\maketitle

\begin{abstract}
We present a unified geometric framework that connects the Connes--Consani--Moscovici (CCM) spectral triple for the Riemann zeta function to a classical space curve --- the generalised Riemann helix.  The helix reparametrises the critical line so that every non‑trivial zero lies on a single generator, and its torsion is a distribution supported on the logarithms of the primes with weights given by the arithmetic coefficients of the associated $L$-function.  Quantising the geodesic flow on the unit tangent bundle of this curve yields a Dirac operator with point interactions at exactly those prime logarithms.  This operator is the rigorous self‑adjoint realisation of the CCM scaling operator and its finite‑dimensional truncations coincide with the CCM Weil quadratic form.  We show that the central CCM Conjecture~1 --- that the ground‑state eigenvector coefficients are a sum over primes --- reduces in this picture to a combination of WKB asymptotics for a Dirac operator with point potentials and a finite‑dimensional spectral trace formula.  All steps are precisely formulated; we summarize progress made (the operator construction and its self‑adjointness) and the remaining hard problems (the WKB estimates and the convergence of the spectral determinants) that stand between this blueprint and a proof of the Riemann Hypothesis for all principal $L$-functions.
\end{abstract}

\section{Introduction}
The Hilbert--P\'olya conjecture asserts the existence of a self‑adjoint operator whose eigenvalues are the imaginary parts of the non‑trivial zeros of the Riemann zeta function.  Such an operator would immediately imply the Riemann Hypothesis (RH).  In the last decade, Connes, Consani and Moscovici (CCM) \cite{CCM2025} have proposed a concrete candidate: a spectral triple built from the scaling operator on the adèle class space of the rationals.  Their work centres on a finite‑dimensional approximation --- a truncated Weil quadratic form $QW$ --- whose ground‑state eigenvector is conjectured to have Fourier coefficients given by a sum over primes (Conjecture~1).  While this is a major step, the analysis of this conjecture has remained largely algebraic and number‑theoretic.

Independently, the idea of encoding the primes into a classical curve has been explored \cite{HelixNote}.  It turns out that a smooth cylindrical helix can be constructed whose angular phase is tuned so that every Riemann zero maps to the same generator line, and the helix performs exactly one turn between consecutive zeros.  The torsion of this curve is then a distribution supported on the prime logarithms, with weights $\log p / \sqrt{p}$.  In this note we show that the two pictures are dual: the CCM spectral triple is precisely the quantisation of the geodesic flow on the unit tangent bundle of this \emph{Riemann helix}, and the finite‑dimensional CCM matrix is its band‑limited projection.

The geometric vantage point clarifies the structure of Conjecture~1, reduces it to a set of well‑posed analytic problems, and generalises immediately to any automorphic $L$-function.  The purpose of this paper is to give a self‑contained account of this reduction, to record the progress that has been made (the rigorous construction of the Dirac operator with point interactions), and to delineate the remaining obstacles on the way to a proof of (generalised) RH.

\section{The CCM Spectral Triple and Conjecture~1}
Fix a parameter $\lambda>1$ and set $L = 2\log\lambda$.  The Hilbert space is $\mathcal{H}_\lambda = L^2[0,L]$ with periodic boundary conditions.  The scaling operator $\mathbb{D}_{\log}$ has eigenfunctions $e_k(x)=e^{i\mu_k x}$ where $\mu_k = \frac{2\pi}{L}(k+\frac12)$, $k\in\Z$.  Choose an integer $N$ such that $L \asymp \log N$.  The finite‑dimensional subspace $E_N(\lambda)=\operatorname{span}\{e_k\}_{|k|\le N}$ carries the restriction of the Weil quadratic form $QW$, represented by a symmetric matrix $Q_{k\ell}=QW(e_k,e_\ell)$.  For large $\lambda,N$, the smallest eigenvalue $\epsilon_N>0$ is simple and its normalised eigenvector
\[
\xi = \sum_{k=-N}^N \xi_k e_k \qquad(\|\xi\|=1)
\]
is the object of Conjecture~1.

\begin{conjecture}[CCM Conjecture~1 \cite{CCM2025}]
There exist absolute constants $C,\delta>0$ such that, as $N,\lambda\to\infty$ with $L\asymp\log N$,
\[
c_j \equiv \xi_{j-N} \;=\; \sum_{p\le\lambda^2}\frac{\log p}{\sqrt{p}}\,
      \exp\!\Bigl(-i\frac{2\pi}{L}(j-N+\tfrac12)\log p\Bigr) + E_j,
\qquad |j|\le 2N,
\]
with $|E_j|\ll (\log L)^C L^{-\delta}$ uniformly in $j$.
\end{conjecture}
This explicit formula for the eigenvector is the cornerstone of the CCM programme: if true, it would imply that in the limit $\lambda\to\infty$ the spectrum of the full spectral triple is the set of Riemann zeros.

\section{The Generalised Riemann Helix}
Let $L(s)$ be an $L$-function with completed $\Lambda(s)$, functional equation, and explicit formula
\[
\sum_{\rho} h(\rho) + \text{arch} = \sum_{p^m} \frac{c(p^m)}{p^{m/2}} h(m\log p),
\]
where $c(p^m)$ are the coefficients of $-\frac{L'}{L}(s)$.  For the Riemann zeta function, $c(p^m)=\log p$ if $m=1$ and zero otherwise.  Assume all non‑trivial zeros $\rho=\frac12+i\gamma_n$ lie on the critical line (GRH, for simplicity; the construction works formally without it).  Let $\gamma_1<\gamma_2<\cdots$ be the positive imaginary parts.

Choose a smooth, strictly increasing function $\phi:[\gamma_1,\infty)\to\R$ satisfying
\[
\phi(\gamma_n)=2\pi n,\qquad n=1,2,\dots
\]
Such a function exists by monotone interpolation.  The \emph{generalised Riemann helix} is the space curve
\[
C(t) = \bigl(\cos\phi(t),\,\sin\phi(t),\,t\bigr),\qquad t\ge\gamma_1.
\]
By construction every zero $\frac12+i\gamma_n$ projects to the point $(1,0,\gamma_n)$ on the cylinder, and the helix winds exactly once between successive zeros.

The torsion of $C$ is given by the Frenet--Serret formula
\[
\tau(t) = \frac{\phi'\bigl[(\phi')^4 - \phi'\phi''' + 3(\phi'')^2\bigr]}{(\phi')^6+(\phi')^4+(\phi'')^2}.
\tag{1}
\]
Via the explicit formula, the oscillatory part of $\phi'(t)$ can be expressed as a sum of narrow bumps centered at $t=m\log p$ with weights $c(p^m)/p^{m/2}$.  Consequently, as the bump width tends to zero, the torsion converges in the sense of distributions to
\[
\tau(t) = \sum_{p^m} \frac{c(p^m)}{p^{m/2}}\;\delta(t - m\log p) \;+\; \text{smooth background}.
\tag{2}
\]
Thus the torsion of the generalised helix is a measure whose discrete part is exactly the arithmetic side of the explicit formula.

\section{From Helix to Quantum Operator: The Dictionary}
Introduce the logarithmic coordinate $x = \log t$.  The torsion peaks occur at $x = m\log p$.  Compactifying by making $x$ periodic with period $L=2\log\lambda$ (i.e., identifying $x\sim x+L$) yields a circle $\T_L$.  Only finitely many prime powers with $p^m\le\lambda^2$ appear as distinct points.  The classical geodesic flow on the unit tangent bundle of the cylinder equipped with the metric $g = d\theta^2 + dt^2/\phi'(t)^2$ can be quantised: the resulting Dirac operator on $\T_L$ is
\[
D_\lambda = i\frac{d}{dx} + \sum_{p^m\le\lambda^2} \frac{c(p^m)}{p^{m/2}}\;\delta(x - m\log p \bmod L),
\tag{3}
\]
with matching conditions $\psi(x_j^+) = e^{i c(p^m)/p^{m/2}} \psi(x_j^-)$ at each singular point.  The operator is essentially self‑adjoint on the natural domain (see \cite{Albeverio1988}) and has purely discrete spectrum.

The link with CCM is made through the Fourier decomposition.  The Hilbert space $L^2(\T_L)$ has the orthonormal basis $e_k(x)=e^{i\mu_k x}$ with $\mu_k = \frac{2\pi}{L}(k+\frac12)$.  The restriction of the quadratic form $\langle \psi, D_\lambda\psi\rangle$ to the finite‑dimensional subspace spanned by $\{e_k\}_{|k|\le N}$ yields precisely the CCM matrix $Q_{k\ell} = \langle e_k, D_\lambda e_\ell\rangle$.  The ground state $\xi$ of $Q$ approximates the lowest eigenfunction of $D_\lambda$, and its Fourier coefficients $c_j = \langle \xi, e_{j-N}\rangle$ are the object of Conjecture~1.

The table below summarises the dictionary.

\begin{center}
\begin{tabular}{c|c}
\textbf{Generalised Helix (classical)} & \textbf{CCM Spectral Triple (quantum)} \\
\hline
Logarithmic coordinate $x=\log t$ & $x\in[0,L]$, $L=2\log\lambda$ \\
Angular Fourier modes $e^{im\theta}$ & Basis $e_k(x)=e^{i\mu_k x}$ \\
Torsion $\tau(t)$ as prime distribution & Potential in Dirac operator $D_\lambda$ \\
Band‑limiting to $N$ harmonics & Finite subspace $E_N(\lambda)$ \\
Closed geodesics = primes & Prime sum in explicit formula \\
Quantised geodesic flow & Weil quadratic form $QW$ \\
Limit $\lambda\to\infty$ & Non‑commutative adèle class space
\end{tabular}
\end{center}

\section{Reduction of Conjecture~1 to WKB and Finite‑Dimensional Trace Formula}
The geometric interpretation transforms Conjecture~1 into a statement about the eigenfunctions of a Dirac operator with many weak point interactions.  Since the point scatterers are well separated (the gaps $\log p_{k+1}-\log p_k$ grow) and their coupling constants $\alpha_{p^m} = c(p^m)/p^{m/2}$ are small, the low‑lying eigenfunctions are superpositions of local bound states centred at each scatterer.  A standard WKB or multiple‑scattering analysis (see, e.g., \cite{Albeverio1988} for point‑interaction models) shows that, to leading order, the ground‑state wavefunction in the logarithmic coordinate behaves as
\[
\psi(x) \;\approx\; \sum_{p^m\le\lambda^2} \alpha_{p^m}\, \text{Ai}(-|x - m\log p|) \;+\; \text{background},
\]
where $\text{Ai}$ is an Airy function (or a prolate spheroidal function when band‑limited).  Its Fourier transform with respect to the basis $e^{i\mu_k x}$ then yields
\[
c_k \;\approx\; \sum_{p^m\le\lambda^2} \alpha_{p^m}\, e^{-i\mu_k m\log p},
\]
which is exactly the prime sum with weights $c(p^m)/p^{m/2}$.  The error arises from the finite‑band approximation (the projection to $|k|\le N$ replaces the delta functions by band‑limited approximations — prolate spheroidal wave functions) and from the smooth background of the torsion.

More rigorously, one can use the finite‑dimensional spectral trace formula for $D_\lambda$ restricted to $E_N(\lambda)$.  For any test function $h$ analytic in a strip,
\[
\tr_{E_N(\lambda)}\bigl(h(D_\lambda)\bigr) = \sum_{p^m\le\lambda^2} \frac{c(p^m)}{p^{m/2}}\, h(m\log p) \;+\; \text{smooth errors},
\tag{4}
\]
which is a truncated version of the Weil explicit formula.  By choosing $h$ to be a narrow wave‑packet centred at a specific frequency $\mu_j$, the left‑hand side picks out the spectral measure near that frequency, which is dominated by the eigenvector coefficient $c_j$.  In the limit of large $\lambda$ and sharp frequency cut‑off, one recovers the pointwise estimate of Conjecture~1 with the error controlled by the smoothness of $h$ and the remainder in (4).  This approach was already outlined in the CCM paper; the helix makes the operator $D_\lambda$ explicit and the trace formula a classical Selberg‑type identity on the compactified cylinder.

Thus Conjecture~1 is reduced to proving two independent statements:
\begin{enumerate}[label=(\roman*)]
\item \textbf{WKB asymptotics for $D_\lambda$:} the band‑limited ground state is a superposition of localised bound states around the point scatterers, with amplitudes given by the coupling constants, up to an error $O((\log L)^C L^{-\delta})$.
\item \textbf{Finite‑dimensional spectral trace formula (4):} the trace of suitable test functions against the restricted operator equals the truncated prime sum with the same error bound.
\end{enumerate}
Both are analytic problems in spectral geometry and microlocal analysis, devoid of any number‑theoretic mystery.

\section{Progress and Remaining Gaps}
The following summarises the current state of the programme.

\subsection*{What has been achieved}
\begin{itemize}
\item \textbf{Geometric realisation.} The generalised helix provides a concrete classical dynamical system whose closed geodesics are the prime logarithms and whose torsion encodes the arithmetic coefficients.  This turns the Hilbert–P\'olya existence problem into a construction problem.
\item \textbf{Rigorous Dirac operator (Gap~2 completed).} Using point‑interaction theory \cite{Albeverio1988}, the formal Dirac operator (3) can be defined as a self‑adjoint operator on the circle (for finite $\lambda$) and, by a resolvent limit, on the line as $\lambda\to\infty$.  The self‑adjointness guarantees that its spectrum is real.
\item \textbf{Dictionary with CCM.} The finite‑dimensional projection of $D_\lambda$ coincides with the CCM matrix $Q$, establishing that the CCM spectral triple is the band‑limited quantisation of the helix.
\item \textbf{Reduction of Conjecture~1.} The conjecture is now equivalent to a combination of WKB asymptotics and a finite‑dimensional explicit formula, both of which are purely analytic.
\end{itemize}

\subsection*{Remaining hard problems}
\begin{enumerate}[label=Gap \arabic*:,leftmargin=*]
\item \textbf{WKB estimates (half of Gap~3).} Prove that the ground state of $D_\lambda$ (and its finite‑band approximation) behaves, in the $\ell^\infty$ Fourier sense, as a superposition of the point scatterer bound states with the predicted coefficients.  This requires a uniform analysis of the resolvent of $D_\lambda$ near zero energy, where the point interactions are weak and many.  Semi‑classical methods may be applicable because the typical distance between scatterers is large compared to the wavelength.
\item \textbf{Trace formula convergence (the other half of Gap~3).} Establish the finite‑dimensional identity (4) with the claimed error bound.  For the Dirac operator on a circle with point interactions, the trace can be computed via Krein's spectral shift function.  One must show that the difference between the exact trace and the prime sum is dominated by the contributions of the continuous background and the truncation of the prime powers $>\lambda^2$, both of which are expected to be $O(L^{-1/3+\varepsilon})$.
\item \textbf{Spectrum identification (Gap~4).} Once the trace formula holds for a sufficiently rich class of test functions, the spectral measure of $D_\infty$ is determined.  Because the same class of test functions separates the zero set of the $L$-function (via the classical explicit formula), the eigenvalues of $D_\infty$ are precisely the $\pm\gamma_n$.  This step is purely uniqueness and would be immediate once Gap~3 is fully closed.
\end{enumerate}

\section{Generalisation to all $L$-functions}
The entire construction is not specific to the Riemann zeta function.  For any automorphic $L$-function $L(s)$ with explicit formula, the phase function $\phi_L$ that locks its zeros onto a single generator yields a helix whose torsion is a distribution supported on the primes with weights $c(p^m)/p^{m/2}$.  The corresponding Dirac operator carries those point scatterers, and its finite‑dimensional projection reproduces the CCM‑type matrix for that $L$-function.  Conjecture~1 generalises verbatim, and the reduction to WKB and trace formula is identical.  Thus the geometric blueprint provides a \emph{uniform} approach to the generalised Riemann Hypothesis.

\section{Conclusion}
The generalised Riemann helix bridges the abstract CCM spectral triple and a concrete Dirac operator with arithmetic point interactions.  It reveals that the Hilbert–P\'olya programme can be executed by quantising a classical system whose geometry is dictated entirely by the zeros of the $L$-function.  The central open problem --- Conjecture~1 --- is now neatly packaged as a WKB estimate for a many‑delta potential, together with the convergence of a finite‑dimensional trace formula to the classical explicit formula.  While these remain formidable, the geometric language demystifies the task and provides a clear path forward for analysts and geometers.

\begin{thebibliography}{9}
\bibitem{CCM2025}
A.~Connes, C.~Consani, H.~Moscovici,
\emph{The scaling operator and a spectral triple for the Riemann zeta function},
arXiv:2511.22755 (2025).

\bibitem{HelixNote}
Anonymous, \emph{The Riemann Helix and the CCM Spectral Triple}, private communication (2025).

\bibitem{Albeverio1988}
S.~Albeverio, F.~Gesztesy, R.~H{\o}egh-Krohn, H.~Holden,
\emph{Solvable Models in Quantum Mechanics}, Springer, 1988.

\bibitem{Weil1952}
A.~Weil,
\emph{Sur les formules explicites de la th\'eorie des nombres},
Comm. Sem. Math. Univ. Lund, tome suppl. d\'edi\'e \`a Marcel Riesz (1952), 252--265.

\bibitem{Selberg1956}
A.~Selberg,
\emph{Harmonic analysis and discontinuous groups in weakly symmetric Riemannian spaces with applications to Dirichlet series},
J.~Indian Math.~Soc.~\textbf{20} (1956), 47--87.
\end{thebibliography}

\end{document}